\chapter{はじめに}
\label{chap:introduction}

\section{不可逆圧縮としての言語}
言語は、人間の経験という連続した海に境界線を引く。アン・サリヴァンが冷たい水をヘレン・ケラーの手のひらに流しながら「W-A-T-E-R」とサインしたその瞬間、分化していない感覚の流れは、名前のついた事物からなる構造的な世界へと変貌した。言語は連続した知覚を離散的かつ伝達可能な単位へと分割する。

しかし、\textbf{言語化は本質的に不可逆的な情報損失を伴う}。生の感情が離散的な語彙へと押し込められる過程で、微細なニュアンスは利用可能な語彙が提供するカテゴリーへと収束し、平坦化されてしまう。「怒り」や「喜び」といった感情語は、言語や文化をまたいで一様に対応するものではない\cite{jackson2019emotion}。同一の言語の中においてさえ、ある瞬間に共存する安堵と孤独の微妙な混在は、「なんとも懐かしい」といった凡庸な表現へと均されてしまうことがある。言語は人間の心と社会をつなぐ唯一の標準化されたインターフェースでありながら、その過程で不可避的に情報を捨象し続けている。


\section{研究目的}

ここで生じる問いは、失われた情報を語彙のボトルネックに到達する\emph{以前}に回収できるか否かである。LLMは感情を何百もの内部層にわたって処理する。モデルがトークンを確定する段階では、細粒度の感情的構造はすでに離散語上の分布へと圧縮されてしまっている。本研究を駆動する仮説は、そのより早い、言語化以前の段階に介入することで、プロンプトエンジニアリング単独では到達し得ない、より精密かつ表現豊かな制御が実現できるというものである。

この仮説は、介入機構の選択という即座の設計上の問いを伴っていた。初期の候補として、感情ラベル付きデータによるファインチューニングやソフトプロンプト最適化が挙げられたが、いずれも推論時に勾配計算を必要とし、解釈が難しい。\textbf{コントラスティブ活性化加算}（CAA）~\cite{caa}はより単純な代替手段であり、対照的な文ペアから差分ベクトルを計算し、それを残差ストリームに加算する手法である。その線形性により、介入と効果の関係が明示的に把握できる。リスクは、単一の線形方向が合成的な感情を捉えるには粗すぎる可能性があることである。それが実際の限界となるかどうかは、本研究が取り組む実証的問いの一つである。

本研究は以下の問いを検討する：
\begin{enumerate}
    \item LLMの内部活性化から信頼性の高い言語化以前の感情表現を抽出することは可能か。すなわち、表層的な語彙的内容から幾何学的に構造化され分離可能な表現は存在するか？
    \item この表現を標的として操作すること、すなわち一つの感情方向を増幅しながら別の方向を抑制することで、モデルの生成出力に一貫した制御可能な変化をもたらすことができるか？
    \item 感情表現は異なる入力コンテキストやモデル条件にわたって一貫した幾何学的構造を共有するか、あるいはそれらは汎化しないコンテキスト固有のアーティファクトに過ぎないか？
\end{enumerate}

三番目の問いは最初の二つよりも後になって生じた。初期実験により、あるコンテキストから抽出されたステアリングベクトルが別のコンテキストに直接適用できないことが示唆され、感情の幾何学的構造（一見安定しているように見えた）が単なるプロンプト依存のショートカット集合に過ぎない可能性が浮かび上がった。この構造が真正なものであるかを検証することが、本プロジェクトの中心的なテーマとなった。
